% =====================================================================
%  MSACL DS301 Deep Learning · Lab 1 · Code Hint Sheet
%  Multiple-choice code options for the two fill-in cells.
%  Style matches quizzes/src/quiz01_foundations.tex (article/fancyhdr).
% =====================================================================
\documentclass[11pt]{article}

\usepackage[letterpaper, margin=0.55in, top=0.55in, bottom=0.4in]{geometry}
\usepackage{amsmath, amssymb}
\usepackage{xcolor}
\usepackage{fancyhdr}
\usepackage{enumitem}

\pagestyle{fancy}
\fancyhf{}
\fancyhead[L]{\small\textbf{MSACL DS301 --- Deep Learning}}
\fancyhead[R]{\small\textbf{Lab 1: Code Hint Sheet}}
\fancyfoot[C]{\thepage}
\renewcommand{\headrulewidth}{0.4pt}

\setlength{\parindent}{0pt}
\setlength{\parskip}{1.5pt}
\newcommand{\opt}[1]{\texttt{#1}}
\setlist[itemize]{leftmargin=2.2em, itemsep=0pt, topsep=0pt, parsep=0pt}

\begin{document}
\small

\begin{center}
{\Large \textbf{Lab 1 --- Code Hint Sheet}}\\[2pt]
{\normalsize Your First Network: the PyTorch training loop}
\end{center}

\vspace{-2pt}
\begin{center}
\fbox{\parbox{0.92\textwidth}{\small
Stuck on a \textbf{YOUR TURN} cell? Each line you need is below with a few options.
Pick one, type it in, and \textbf{run the cell} --- the built-in \texttt{assert} checks will tell you
whether it is right. Only the two marked cells need editing; run everything else as-is.
}}
\end{center}

\vspace{2pt}\hrule\vspace{4pt}

\textbf{Step 4 --- the training loop.}\; {\small Write these five lines \emph{in order} inside the
\texttt{for epoch} loop (Lecture 2: forward $\to$ loss $\to$ zero $\to$ backward $\to$ step).}

\textbf{1. Forward --- make a guess.}
\begin{itemize}
  \item[(a)] \opt{pred = model(X\_train)}
  \item[(b)] \opt{pred = loss\_fn(X\_train)}
  \item[(c)] \opt{pred = model.guess(y\_train)}
\end{itemize}

\textbf{2. Loss --- how wrong?}
\begin{itemize}
  \item[(a)] \opt{loss = loss\_fn(pred, y\_train)}
  \item[(b)] \opt{loss = model(pred, y\_train)}
  \item[(c)] \opt{loss = loss\_fn(X\_train, y\_train)}
\end{itemize}

\textbf{3. Zero --- clear last step's gradients.}
\begin{itemize}
  \item[(a)] \opt{optimizer.zero\_grad()}
  \item[(b)] \opt{optimizer.reset()}
  \item[(c)] \opt{model.zero\_grad(loss)}
\end{itemize}

\textbf{4. Backward --- backprop the blame.}
\begin{itemize}
  \item[(a)] \opt{loss.backward()}
  \item[(b)] \opt{optimizer.backward()}
  \item[(c)] \opt{model.backward(loss)}
\end{itemize}

\textbf{5. Step --- nudge every weight.}
\begin{itemize}
  \item[(a)] \opt{optimizer.step()}
  \item[(b)] \opt{optimizer.update()}
  \item[(c)] \opt{loss.step()}
\end{itemize}

\vspace{3pt}\hrule\vspace{4pt}

\textbf{Step 6 --- the learning rate: pick a too-big value, then read the curves.}\; {\small You are
given \opt{lr\_ok = 0.01}. Choose a rate that is clearly \emph{too big} so the loss climbs, run both,
then decide which one diverged.}

\textbf{1. A learning rate that is clearly TOO BIG.}
\begin{itemize}
  \item[(a)] \opt{lr\_too\_big = 5.0}\quad{\small (huge step $\Rightarrow$ overshoots, loss rises)}
  \item[(b)] \opt{lr\_too\_big = 0.001}\quad{\small (smaller than \opt{lr\_ok} --- trains fine, doesn't diverge)}
  \item[(c)] \opt{lr\_too\_big = 0.01}\quad{\small (equal to \opt{lr\_ok} --- no contrast at all)}
\end{itemize}

\textbf{2. Which run diverged? (read it off the two-curve plot).}
\begin{itemize}
  \item[(a)] \opt{which\_diverged = "too\_big"}\quad{\small (its loss went UP, not down)}
  \item[(b)] \opt{which\_diverged = "ok"}\quad{\small (no --- the sensible rate fell smoothly)}
  \item[(c)] \opt{which\_diverged = "both"}\quad{\small (no --- only the too-big curve rose)}
\end{itemize}

\vspace{3pt}\hrule\vspace{4pt}

\textbf{Step 7 --- evaluate on the test set.}\; {\small Predict without learning, threshold at 0.5,
then measure the fraction correct.}

\textbf{1. Predict without gradients.}
\begin{itemize}
  \item[(a)] \opt{with torch.no\_grad():}\quad{\small (then, indented)}\quad\opt{probs = model(X\_test)}
  \item[(b)] \opt{with torch.grad():}\quad{\small (then, indented)}\quad\opt{probs = model(X\_test)}
  \item[(c)] \opt{probs = model.no\_grad(X\_test)}
\end{itemize}

\textbf{2. Turn probabilities into 0/1 predictions.}
\begin{itemize}
  \item[(a)] \opt{predicted = (probs > 0.5).float()}
  \item[(b)] \opt{predicted = probs.round(threshold=0.5)}
  \item[(c)] \opt{predicted = probs.argmax()}
\end{itemize}

\textbf{3. Accuracy --- fraction that match the label.}
\begin{itemize}
  \item[(a)] \opt{test\_acc = (predicted == y\_test).float().mean().item()}
  \item[(b)] \opt{test\_acc = (predicted - y\_test).mean()}
  \item[(c)] \opt{test\_acc = accuracy(predicted, y\_test)}
\end{itemize}

\vspace{4pt}\hrule\vspace{4pt}

\textbf{Step 8 --- did it overfit? compute the training accuracy the same way.}\; {\small Mirror the
Step 7 evaluation, but on \opt{X\_train}/\opt{y\_train}, then compare to \opt{test\_acc}.}

\textbf{1. Training accuracy (same recipe as Step 7, on the training data).}
\begin{itemize}
  \item[(a)] {\small\opt{with torch.no\_grad():}\; \opt{probs\_train = model(X\_train)}\; then}\\
        \hspace*{1.2em}\opt{predicted\_train = (probs\_train > 0.5).float()}\\
        \hspace*{1.2em}\opt{train\_acc = (predicted\_train == y\_train).float().mean().item()}
  \item[(b)] \opt{train\_acc = (model(X\_test) == y\_test).float().mean()}\quad{\small (that's the \emph{test} set, not train)}
  \item[(c)] \opt{train\_acc = loss\_history[-1]}\quad{\small (that's a loss, not an accuracy)}
\end{itemize}

\textbf{2. Overfitting means the model scores HIGHER on data it has seen. Fill the comparison.}
\begin{itemize}
  \item[(a)] \opt{overfitting = train\_acc > test\_acc}\quad{\small (train beats test $\Rightarrow$ overfit)}
  \item[(b)] \opt{overfitting = train\_acc < test\_acc}\quad{\small (backwards --- and the assert will fail)}
  \item[(c)] \opt{overfitting = train\_acc == test\_acc}\quad{\small (they're not equal here)}
\end{itemize}

\vspace{4pt}\hrule\vspace{3pt}
{\small\itshape Answers self-check in the notebook: if the \texttt{assert} cells print their success
message, you chose right. The loss should fall steeply and then flatten; the too-big learning rate
makes the loss \emph{rise} instead; test accuracy should land around \textbf{0.72--0.77} (coin-flip
baseline is 0.50); and training accuracy should be clearly higher than test accuracy (a train/test
gap of roughly \textbf{0.20--0.25}) --- that gap is overfitting.}

\end{document}
